% ============================================================
% Hesam Salmabadi — Resume
% Target: Geospatial Data Scientist, Innovate! Inc. / USGS EROS
% Compiled with: xelatex resume_innovate_geospatial_ds.tex
% ============================================================

\documentclass[10pt, a4paper]{article}

% --- Packages ---
\usepackage{geometry}
\geometry{top=1.6cm, bottom=1.6cm, left=2cm, right=2cm}

\usepackage{fontspec}
\setmainfont{Liberation Sans}

\usepackage{titlesec}
\usepackage{enumitem}
\usepackage{hyperref}
\usepackage{parskip}
\usepackage{microtype}
\usepackage{xcolor}

% --- Colours ---
\definecolor{headingblue}{HTML}{1A3A5C}
\definecolor{linkblue}{HTML}{1A5276}
\definecolor{lightgray}{HTML}{666666}

% --- Hyperref ---
\hypersetup{
    colorlinks=true,
    urlcolor=linkblue,
    linkcolor=linkblue,
    pdfauthor={Hesam Salmabadi},
    pdftitle={Hesam Salmabadi — Resume},
}

% --- Section formatting ---
\titleformat{\section}
  {\large\bfseries\color{headingblue}}
  {}{0em}{}
  [\vspace{-4pt}\textcolor{headingblue}{\rule{\linewidth}{0.6pt}}\vspace{2pt}]

\titlespacing{\section}{0pt}{10pt}{4pt}

\pagestyle{empty}

\setlist[itemize]{noitemsep, topsep=2pt, leftmargin=1.4em}
\setlist[itemize,1]{label=\textbullet}

% --- Custom commands ---
\newcommand{\dategray}[1]{\textcolor{lightgray}{\small #1}}
\newcommand{\jobtitle}[4]{%
  \noindent\textbf{#1} \hfill \dategray{#3} \\
  \textit{#2} \hfill \dategray{#4} \\[-4pt]
}
\newcommand{\publication}[1]{\noindent #1 \vspace{3pt}}

% ============================================================
\begin{document}
% ============================================================

% --- HEADER ---
\begin{center}
  {\LARGE\bfseries\color{headingblue} Hesam Salmabadi}\\[4pt]
  {\small
    Montreal, QC, Canada \quad|\quad
    \href{mailto:hesam.salmabadi@gmail.com}{hesam.salmabadi@gmail.com} \quad|\quad
    +1 514 430 2051 \\[2pt]
    \href{https://hesam-salmabadi.github.io}{hesam-salmabadi.github.io} \quad|\quad
    \href{https://github.com/hesam-salmabadi}{GitHub} \quad|\quad
    \href{https://www.linkedin.com/in/hesam-salmabadi/}{LinkedIn} \quad|\quad
    \href{https://scholar.google.com/citations?user=QAvXX3cAAAAJ}{Google Scholar}
  }
\end{center}

\vspace{2pt}

% ============================================================
\section{Profile}
% ============================================================

Geospatial data scientist with nine peer-reviewed publications and direct experience
designing national-scale land use/land cover (LULC) modeling and projection
frameworks. Currently at Agriculture and Agri-Food Canada building ML-based
suitability models and Bayesian demand forecasting pipelines for continental-scale
satellite datasets. Proficient in Python geospatial stack (xarray, dask, GDAL,
scikit-learn), spatial cross-validation, and multi-source satellite data integration.
Contributing author, \textit{Canada's Changing Climate Report 2026}.

% ============================================================
\section{Skills}
% ============================================================

\noindent\textbf{Programming \& Data Engineering:} Python (xarray, dask, numpy,
pandas, rasterio, geopandas, scikit-learn, PyMC, matplotlib, cartopy, h5py),
GDAL/OGR (Python \& CLI), Bash, Git/GitHub, LaTeX, MATLAB

\vspace{3pt}
\noindent\textbf{Machine Learning \& Statistics:} Logistic regression (L1/L2),
random forests, XGBoost, SVM, Bayesian hierarchical modelling (PyMC / NUTS),
spatial block cross-validation, imbalanced classification, Monte Carlo uncertainty
propagation, Mann-Kendall trend analysis

\vspace{3pt}
\noindent\textbf{Remote Sensing \& Earth Observation:} Passive microwave (SMAP,
AMSR2), multispectral satellite time-series (Landsat, MODIS), atmospheric transport
modelling (HYSPLIT), ERA5 climate reanalysis, multi-sensor data fusion

\vspace{3pt}
\noindent\textbf{GIS \& Spatial Analysis:} QGIS, ArcGIS/ArcMap, ArcPy, geostatistical
interpolation (IDW, kriging), kernel density estimation

\vspace{3pt}
\noindent\textbf{Data Formats:} Zarr, NetCDF, HDF4/5, GeoTIFF, GRIB, Parquet,
Shapefile --- national-scale out-of-core processing pipelines

% ============================================================
\section{Experience}
% ============================================================

\jobtitle{Geospatial Data Scientist}
         {Agriculture and Agri-Food Canada (AAFC), Science \& Technology Branch}
         {May 2025 -- Present}
         {Ottawa, ON, Canada (Remote)}
\begin{itemize}
  \item Designed and implemented a national-scale LULC modeling and projection
        framework for Canada (5 classes: Agriculture, Urban, Natural Vegetation,
        Wetland, Water) at 100\,m resolution, integrating multi-temporal satellite
        time series and 16+ environmental covariates --- scope and objectives
        directly parallel to USGS EROS land cover modeling projects
  \item Developed binary logistic regression classifiers for agricultural, urban,
        and natural vegetation suitability mapping using quadratic feature
        expansion; implemented 75$\times$75\,km spatial blocking cross-validation
        (Stratified Group K-Fold) to prevent data leakage in spatially
        autocorrelated satellite datasets (ROC-AUC 0.957, Recall 0.81 on
        geographically held-out blocks)
  \item Built Bayesian demand forecasting models (PyMC, NUTS sampler) for
        provincial LULC trajectories across three independent data sources
        (satellite, administrative, census); reported posterior medians and
        94\% HDI to quantify projection uncertainty
  \item Engineered out-of-core national raster pipelines (Python: xarray, dask,
        rasterio, GDAL) ingesting and harmonising multi-temporal LULC time series
        into a cloud-optimised Zarr database (\textit{canada\_lulc\_100m\_v4.zarr})
        covering all of Canada south of 60°N
  \item Cross-validated satellite-derived LULC estimates against Statistics Canada
        Census across all 10 southern provinces; diagnosed three distinct
        divergence regimes with physical explanations (Pearson $r = 0.988$,
        normalised by province land area)
\end{itemize}

\vspace{6pt}

\jobtitle{PhD Researcher}
         {Université du Québec à Trois-Rivières (UQTR) \& Centre d'Études Nordiques, Université Laval}
         {May 2021 -- May 2026}
         {Trois-Rivières, QC, Canada}
\begin{itemize}
  \item Developed a non-binary soil freeze–thaw detection framework using
        multi-frequency passive microwave satellite data; validated across
        87 monitoring sites spanning boreal, prairie, and tundra ecozones
        (match rates 78.8–98.2\%, Kappa 0.44–0.87)
  \item Designed and trained probabilistic spatial classifiers (logistic
        regression, Bayesian hierarchical models) with rigorous spatial
        block cross-validation and imbalanced-class evaluation (PRAUC, EER)
  \item Processed and analysed multi-year satellite archives (SMAP, AMSR2,
        ERA5) using Python (xarray, dask, Zarr, HDF5/NetCDF) at continental
        scale; built batch-and-append pipelines and ML-ready Parquet exports
  \item Funded by FRQNT Doctoral Research Scholarship (\$75,000 CAD,
        2023–2026) and UQTR Bourse Universalis Causa
\end{itemize}

\vspace{6pt}

\jobtitle{Geospatial Data Analyst --- Contributing Author}
         {Environment and Climate Change Canada (ECCC)}
         {Apr 2024 -- Oct 2024}
         {Remote}
\begin{itemize}
  \item Contributing author, \textit{Canada's Changing Climate Report 2026},
        Chapter 6: Cryosphere Changes --- delivered geospatial analysis and
        publication-quality maps for a national federal assessment report
  \item Built a parallel ETL pipeline (ProcessPoolExecutor) to ingest and
        process 40+ years of NSIDC satellite freeze–thaw records (HDF5);
        applied ArcPy automation for cartographic delivery
\end{itemize}

\vspace{6pt}

\jobtitle{Environmental Remote Sensing Scientist}
         {Research Institute of Cultural Heritage \& Tourism (RICHT)}
         {Oct 2019 -- Jan 2021}
         {Tehran, Iran}
\begin{itemize}
  \item Investigated long-range atmospheric transport of dust from Lake
        Urmia's desiccating playa using HYSPLIT forward trajectories and
        multi-decadal MODIS/Landsat satellite time-series analysis
\end{itemize}

\vspace{6pt}

\jobtitle{Field Research Technician}
         {NASA Jet Propulsion Laboratory (JPL) --- SMAPVEX22-Boreal Campaign}
         {Aug 2022}
         {Saskatchewan, Canada}
\begin{itemize}
  \item Contributed soil sampling, sensor deployment, and LiDAR data
        collection across 37 sites in support of SMAP satellite soil moisture
        validation; field data published in Berg et al.\ (2025),
        \textit{IEEE JSTARS}
\end{itemize}

% ============================================================
\section{Education}
% ============================================================

\jobtitle{Doctor of Philosophy --- Environmental Sciences}
         {Université du Québec à Trois-Rivières (UQTR)}
         {May 2021 -- May 2026}
         {Trois-Rivières, QC, Canada}

\noindent\textit{Thesis:} Non-Binary Monitoring of Seasonally Frozen Ground:
From In Situ Sensors to Continental Passive Microwave Retrievals \quad
\textit{Supervisor:} Prof.\ Alexandre Roy

\vspace{6pt}

\jobtitle{MSc --- Civil Engineering (Environmental Engineering)}
         {Iran University of Science and Technology (IUST)}
         {Sep 2015 -- Jul 2018}
         {Tehran, Iran}

\noindent\textit{GPA:} 18.33/20 \quad First Ranked Student \quad
\textit{Thesis:} Determining the Effect of Wetland and Lake Desiccation in Iran
on the Air Quality of Potentially Affected Regions

\vspace{6pt}

\jobtitle{BSc --- Civil Engineering}
         {University of Birjand}
         {Sep 2010 -- Sep 2014}
         {Birjand, Iran}

% ============================================================
\section{Selected Publications}
% ============================================================

\noindent\textcolor{lightgray}{\small
  14 publications total (9 first-authored) \,|\,
  Citations: 181 \,|\, h-index: 6 \,|\, i10-index: 5 \quad
  \href{https://scholar.google.com/citations?user=QAvXX3cAAAAJ}{[Google Scholar]}
}

\vspace{4pt}

\publication{\textbf{Salmabadi H.}, Berg A., Roy A. (2026, preprint).
Non-Binary Retrievals of Soil Freeze–Thaw Dynamics Using Multi-Frequency
Passive Microwave Data. \textit{SSRN}, under review.
\href{https://ssrn.com/abstract=6626104}{ssrn.com/abstract=6626104}}

\publication{\textbf{Salmabadi H.}, Pardo Lara R., Berg A., et al.\ (2026).
In situ monitoring of seasonally frozen ground using soil freezing characteristic
curve in permittivity–temperature space. \textit{The Cryosphere}, 20, 1635–1654.
\href{https://doi.org/10.5194/tc-20-1635-2026}{doi:10.5194/tc-20-1635-2026}}

\publication{\textbf{Salmabadi H.}, Saeedi M., Notaro M., Roy A. (2025).
Dust transport pathways from the Mesopotamian Marshes.
\textit{Aeolian Research}, 73, 100975.
\href{https://doi.org/10.1016/j.aeolia.2025.100975}{doi:10.1016/j.aeolia.2025.100975}}

\publication{Berg A.A., \ldots, \textbf{Salmabadi H.}, \ldots\ (2025).
Soil Moisture Active/Passive Validation Experiment Within the Canadian Boreal
Forest (SMAPVEX22-Boreal). \textit{IEEE J. Sel. Topics Appl. Earth Obs.
Remote Sens.}, 18.
\href{https://doi.org/10.1109/JSTARS.2025.3564195}{doi:10.1109/JSTARS.2025.3564195}}

% ============================================================
\section{Awards \& Funding}
% ============================================================

\noindent\textbf{FRQNT Doctoral Research Scholarship} \hfill \dategray{2023--2026}\\
Fonds de recherche du Québec – Nature et technologies \quad \$75,000 CAD

\vspace{4pt}
\noindent\textbf{First Ranked Student --- MSc Program} \hfill \dategray{2018}\\
Iran University of Science and Technology \quad GPA: 18.33/20

\vspace{4pt}
\noindent\textbf{Full Tuition Scholarship (BSc \& MSc)} \hfill \dategray{2010--2018}\\
Ministry of Science, Research and Technology of Iran --- national competitive examination

% ============================================================
\section{Languages}
% ============================================================

Persian (Native) \quad|\quad English (Fluent) \quad|\quad French (Functional)

% ============================================================
\end{document}
